\section{Methodology}

\subsection{Problem Formulation}

\textbf{Input:} A state with $N$ census tracts, each with population $p_i$ and geometry $G_i$. Target: $K$ congressional districts with equal population $\bar{p} = (\sum_i p_i) / K$.

\textbf{Constraints:}
\begin{enumerate}
    \item \textbf{Population balance}: $|p(D_j) - \bar{p}| \leq 0.005 \bar{p}$ for all districts $D_j$ (within 0.5\%)
    \item \textbf{Contiguity}: Each district forms a connected component in the adjacency graph
\end{enumerate}

\textbf{Objective:} Minimize total district perimeters $\sum_{j=1}^{K} P(D_j)$, equivalently maximizing Polsby-Popper compactness since area is held constant.

\textbf{Approach:} Model the state as a weighted graph $G = (V, E, w)$ where vertices $V$ are census tracts, edges $E$ connect adjacent tracts, and weights $w_{ij}$ are boundary lengths. Apply recursive METIS bisection with edge weights to minimize weighted edge cuts, which directly corresponds to minimizing total perimeter.

\subsection{Graph Construction}

\subsubsection{Adjacency Detection}

We use \textbf{Queen contiguity}: tracts $i$ and $j$ are adjacent if their boundaries share an edge or vertex (not merely nearest neighbors). Implemented via spatial indexing (R-tree) with intersection tests:

\begin{algorithmic}[1]
\Function{DetectAdjacency}{$\{G_1, \ldots, G_N\}$}
    \State $E \gets \emptyset$
    \State Build R-tree index from tract bounding boxes
    \For{each tract $i$}
        \State candidates $\gets$ R-tree query($G_i$.bounds)
        \For{each candidate $j > i$}
            \If{$G_i$.intersects($G_j$)}
                \State $E \gets E \cup \{(i, j)\}$
            \EndIf
        \EndFor
    \EndFor
    \State \Return $E$
\EndFunction
\end{algorithmic}

Complexity: $O(N \log N)$ for R-tree construction, $O(Nd)$ for intersection tests with average degree $d \approx 6$.

\subsubsection{Boundary Length Computation}

For each edge $(i,j) \in E$, compute the geometric intersection of boundaries:

\begin{equation}
    I_{ij} = \partial G_i \cap \partial G_j
\end{equation}

where $\partial G_i$ denotes the boundary of geometry $G_i$. The intersection type determines the weight:

\begin{equation}
    w_{ij} = \begin{cases}
        \text{length}(I_{ij}) & \text{if } I_{ij} \text{ is LineString/MultiLineString} \\
        0.1 \text{ meters} & \text{if } I_{ij} \text{ is Point (corner contact)} \\
        w_{\text{bridge}} & \text{if } I_{ij} = \emptyset \text{ (water crossing)}
    \end{cases}
\end{equation}

\textbf{Geometric operations} use Shapely~\cite{shapely} with state-specific projections (e.g., California Albers EPSG:3310, Texas Conic EPSG:3083) for accurate length measurement. Point contacts receive near-zero weight (0.1m) to permit adjacency without contributing to perimeter.

\subsubsection{County-Based Bridge Connections}

For states with water barriers (30+ states including island states and those with major lakes/rivers), adjacency graphs may be disconnected. We add \textbf{bridge edges} to connect components:

\begin{algorithmic}[1]
\Function{AddBridges}{$V, E, w$}
    \State components $\gets$ ConnectedComponents($V, E$)
    \For{each disconnected component $C$}
        \For{each tract $i \in C$}
            \State county($i$) $\gets$ first 5 digits of 11-digit GEOID
            \State target $\gets$ nearest tract $j$ with county($j$) = county($i$) and $j \notin C$
            \If{target exists}
                \State $E \gets E \cup \{(i, \text{target})\}$
                \State $w_{i,\text{target}} \gets$ median($\{w_{kl} \mid (k,l) \in E_{\text{land}}\}$)
            \EndIf
        \EndFor
    \EndFor
\EndFunction
\end{algorithmic}

Bridge edges use \textit{median land boundary length} as weight, providing adaptive scaling: small states have shorter typical boundaries ($\sim$1 km), large states longer ($\sim$5-10 km). This balances permitting necessary water crossings while discouraging excessive splits.

\textbf{County restriction} ensures geographic sensibility: Alaska's Aleutian Islands (~1700 km apart) connect within their shared county rather than to mainland tracts, and Hawaii's islands remain separated by county. This prevents unrealistic long-distance connections while ensuring all tracts are reachable.

\subsection{METIS Integration}

\subsubsection{CSR Format with Edge Weights}

METIS requires Compressed Sparse Row (CSR) format. We use format code \texttt{011}: edge weights + vertex weights (populations).

\begin{verbatim}
<n_vertices> <n_edges> 011 1
<vweight_1> <neighbor_1> <eweight_1> <neighbor_2> <eweight_2> ...
<vweight_2> <neighbor_1> <eweight_1> ...
...
\end{verbatim}

Boundary lengths (continuous meters) are scaled to \textbf{integer centimeters}:

\begin{equation}
    w_{ij}^{\text{METIS}} = \max(1, \lfloor w_{ij} \times 100 \rfloor)
\end{equation}

This provides 1cm precision (adequate for tract-level data) while avoiding METIS numerical issues with floating-point weights. METIS uses 1-based indexing; neighbor indices are incremented during file writing.

\subsubsection{METIS Configuration}

We invoke METIS via command line:

\begin{verbatim}
gpmetis -contig -minconn -ufactor=1.005 -niter=100 graph.txt 2
\end{verbatim}

\begin{itemize}
    \item \texttt{-contig}: Enforce partition contiguity (required for redistricting)
    \item \texttt{-minconn}: Minimize subgraph connectivity (reduces perimeter)
    \item \texttt{-ufactor=1.005}: Allow 0.5\% population imbalance (constitutional threshold)
    \item \texttt{-niter=100}: Refinement iterations for quality
\end{itemize}

METIS objective function with edge weights:

\begin{equation}
    \min_{P \in \{0,1\}^N} \sum_{(i,j) \in E} w_{ij} \cdot \mathbb{1}[P_i \neq P_j]
\end{equation}

subject to population balance and contiguity. Minimizing weighted edge cuts directly minimizes total perimeter.

\subsection{Recursive Bisection Algorithm}

Edge-weighted recursive bisection proceeds top-down:

\begin{algorithm}
\caption{Edge-Weighted Recursive Bisection}
\begin{algorithmic}[1]
\Function{RecursiveBisect}{$V, E, w, K$}
    \If{$K = 1$}
        \State \Return $\{V\}$ \Comment{Base case: single district}
    \EndIf

    \State $k_1 \gets \lfloor K/2 \rfloor$, $k_2 \gets K - k_1$ \Comment{Target sizes}
    \State $(V_1, V_2) \gets$ \Call{METIS-Bisect}{$V, E, w, k_1, k_2$}

    \State Extract subgraphs $G_1 = (V_1, E_1, w_1)$ and $G_2 = (V_2, E_2, w_2)$
    \State $D_1 \gets$ \Call{RecursiveBisect}{$V_1, E_1, w_1, k_1$}
    \State $D_2 \gets$ \Call{RecursiveBisect}{$V_2, E_2, w_2, k_2$}

    \State \Return $D_1 \cup D_2$
\EndFunction
\end{algorithmic}
\end{algorithm}

\textbf{Subgraph extraction} requires index remapping. For region with tracts $\{v_1, \ldots, v_m\}$, create mappings:
\begin{itemize}
    \item $L(i) = v_i$ (local-to-global index)
    \item $G^{-1}(v_i) = i$ (global-to-local index)
\end{itemize}

Extract subgraph edges and weights:
\begin{equation}
    E' = \{(G^{-1}(i), G^{-1}(j)) \mid i,j \in \{v_1, \ldots, v_m\}, (i,j) \in E\}
\end{equation}
\begin{equation}
    w' = \{w_{ij} \mid i,j \in \{v_1, \ldots, v_m\}, (i,j) \in E\}
\end{equation}

\textbf{Unequal splits} handle odd district counts. For $K = 7$, splits are $(4,3), (2,2), (2,1), (1,1), (1,1), (1,1), (1,1)$. METIS supports unequal target partition sizes through vertex weight ratios.

\subsection{Complexity Analysis}

\begin{itemize}
    \item \textbf{Preprocessing}: $O(N \log N)$ for R-tree + $O(Nd)$ for intersections = $O(Nd)$ overall ($d \approx 6$)
    \item \textbf{METIS calls}: $O(N)$ per bisection, $O(\log K)$ levels = $O(N \log K)$ total
    \item \textbf{Overall}: $O(Nd + N \log K)$ dominated by $O(N \log K)$ for large $K$
\end{itemize}

\textbf{Runtime overhead}: Edge-weighted mode adds 10-50\% to baseline (larger METIS arrays, more refinement). Preprocessing (boundary computation) requires 10-30 seconds per state, cached for reuse.

\textbf{Memory}: Edge weights require $\sim$100 KB per 1,000 edges (8-byte floats + indices). Full state graphs: 1-5 MB.

\subsection{Implementation Details}

\textbf{Software stack}: Python 3.13, GeoPandas for geometry, METIS 5.1.0, Shapely 2.0 for spatial operations. Adjacency and weights cached in Parquet format for reproducibility.

\textbf{Parallelization}: States processed independently with 4-6 worker processes (bounded by memory for large states).

\textbf{Validation}: Post-processing verifies population balance ($\pm 0.5\%$), contiguity (BFS traversal), and computes Polsby-Popper scores for all districts.

Open-source implementation available at \url{https://github.com/TODO}.
